\section*{Chapitre \ref{ch:g12:exp} --- Exponentielle et logarithme}

\begin{solution}{exo:g12:exp:1}
$\dfrac{\eu^{3x}\,\eu^{-x+1}}{\eu^{x}} = \eu^{3x - x + 1 - x} = \eu^{x+1}$.

$\ln\!\left(\dfrac{\eu^2\sqrt{\eu}}{\eu^{-3}}\right)
= 2 + \tfrac12 + 3 = \tfrac{11}{2}$.
\end{solution}

\begin{solution}{exo:g12:exp:2}
\emph{(a)} Poser $X = \eu^x > 0$~: $X^2 - 3X + 2 = 0$ donne $X = 1$ ou
$X = 2$, tous deux positifs, donc $x = 0$ ou $x = \ln 2$.

\emph{(b)} \omterm{def:g11:func:function}{Ensemble de définition}~: $x - 1 > 0$ et $x + 2 > 0$, donc
$x > 1$. L'\omterm{def:g10:algebra:equation}{équation} devient $\ln\bigl((x-1)(x+2)\bigr) = \ln 4$, d'où
$(x-1)(x+2) = 4$, \emph{c'est-à-dire} $x^2 + x - 6 = 0$, donc $x = 2$ ou
$x = -3$. Seul $x = 2$ est dans l'\omterm{def:g11:func:function}{ensemble de définition}~: $S = \{2\}$.
\end{solution}

\begin{solution}{exo:g12:exp:3}
Diviser par $\eu^x$~:
$\dfrac{1 - x^2\eu^{-x}}{1 + \eu^{-x}} \to \dfrac{1-0}{1+0} = 1$, en
utilisant $x^2 \eu^{-x} \to 0$ (\cref{thm:g12:exp:growth}).

$\dfrac{\ln x}{\sqrt x} \to 0$ par la \cref{prop:g12:exp:lnanalytic} (cas
$n = 2$).

$x^2 \ln x = x \cdot (x \ln x) \to 0 \times 0 = 0$.

$x - \ln x = x\left(1 - \frac{\ln x}{x}\right) \to +\infty$ car
$\frac{\ln x}{x} \to 0$.
\end{solution}

\begin{solution}{exo:g12:exp:4}
$f'(x) = \eu^{-x} - x\eu^{-x} = (1 - x)\eu^{-x}$, du signe de $1 - x$~:
$f$ croît sur $\intoc{-\infty}{1}$, décroît sur $\intco{1}{+\infty}$,
avec un \omterm{def:g10:functions:extrema}{maximum} global $f(1) = \eu^{-1}$.

Limites~: lorsque $x \to +\infty$, $f(x) = \frac{x}{\eu^x} \to 0$
(\cref{thm:g12:exp:growth})~; lorsque $x \to -\infty$, $x \to -\infty$ et
$\eu^{-x} \to +\infty$, donc $f(x) \to -\infty$.

$f''(x) = -\eu^{-x} - (1-x)\eu^{-x} = (x - 2)\eu^{-x}$, qui change de
signe en $x = 2$~: \omterm{def:g12:deriv:inflection}{point d'inflexion} en $\bigl(2,\, 2\eu^{-2}\bigr)$.
\end{solution}

\begin{solution}{exo:g12:exp:5}
Soit $g(x) = x - \ln(1+x)$ sur $\intoo{-1}{+\infty}$. Alors
$g'(x) = 1 - \frac{1}{1+x} = \frac{x}{1+x}$, négatif sur $\intoo{-1}{0}$
et positif sur $\intoo{0}{+\infty}$~: $g$ atteint son \omterm{def:g10:functions:extrema}{minimum} $g(0) = 0$,
donc $g \geq 0$ avec égalité seulement en $0$. D'où $\ln(1+x) \leq x$.

Appliquer avec $x = \frac1n$~:
$n\ln\left(1 + \frac1n\right) \leq 1$, donc
$\left(1+\frac1n\right)^n = \eu^{\,n\ln(1+1/n)} \leq \eu$.
Appliquer avec $x = -\frac{1}{n+1} > -1$~:
$\ln\left(1 - \frac{1}{n+1}\right) \leq -\frac{1}{n+1}$, donc
$-(n+1)\ln\left(1 - \frac{1}{n+1}\right) \geq 1$ et
$\left(1 - \frac{1}{n+1}\right)^{-(n+1)} \geq \eu$.
\end{solution}

\begin{solution}{exo:g12:exp:6}
\emph{1.} Chaque année multiplie le capital par $1.03$~:
$C_n = C_0 (1.03)^n$.

\emph{2.} $C_n \geq 2C_0 \iff (1.03)^n \geq 2 \iff n \ln 1.03 \geq \ln 2
\iff n \geq \frac{\ln 2}{\ln 1.03} \approx 23.45$~: le capital double
après $24$ années.

\emph{3.} $\frac{70}{3} \approx 23.3$, proche de l'exact
$\frac{\ln 2}{\ln 1.03}$. La règle fonctionne car
$\ln(1 + r) \approx r$ pour $r$ petit, donc
$\frac{\ln 2}{\ln(1+r)} \approx \frac{0.693}{r} \approx \frac{70}{100r}$.
\end{solution}

\begin{solution}{exo:g12:exp:7}
$\eu^{2x} > \eu^{x+1} \iff 2x > x + 1 \iff x > 1$ (l'exponentielle est
strictement \omterm{def:g12:seq:monotonic}{croissante})~: $S = \intoo{1}{+\infty}$.

\omterm{def:g11:func:function}{Ensemble de définition} de la seconde inégalité~: $x^2 - 1 > 0$ et
$x + 5 > 0$, \emph{c'est-à-dire}
$x \in \intoo{-5}{-1} \cup \intoo{1}{+\infty}$. Sur ce domaine, $\ln$
étant strictement \omterm{def:g12:seq:monotonic}{croissante},
\[
\ln(x^2-1) \leq \ln(x+5) \iff x^2 - 1 \leq x + 5 \iff x^2 - x - 6 \leq 0
\iff x \in \intcc{-2}{3}.
\]
En intersectant avec le domaine~: $S = \intco{-2}{-1} \cup \intoc{1}{3}$.
\end{solution}

\begin{solution}{exo:g12:exp:8}
\emph{1.} $f'(x) = \dfrac{\eu^x x - \eu^x}{x^2} = \dfrac{(x-1)\eu^x}{x^2}$,
négatif sur $\intoo{0}{1}$, positif sur $\intoo{1}{+\infty}$~: \omterm{def:g10:functions:extrema}{minimum}
$f(1) = \eu$. Limites~: $f \to +\infty$ en $0^+$ (numérateur $\to 1$,
dénominateur $\to 0^+$) et en $+\infty$ (\cref{thm:g12:exp:growth}).

\emph{2.} Pour $x > 0$, $\eu^x = kx \iff f(x) = k$. D'après le \omterm{def:g10:functions:table}{tableau de
variations} ($+\infty \searrow \eu \nearrow +\infty$, \omterm{def:g12:limcont:continuity}{continue} et
strictement \omterm{def:g11:func:monotone}{monotone} de chaque côté)~: aucune solution pour $k < \eu$~;
exactement une ($x = 1$) pour $k = \eu$~; exactement deux pour $k > \eu$
(une dans $\intoo{0}{1}$, une dans $\intoo{1}{+\infty}$, par le théorème
de la bijection sur chaque \omterm{def:g10:numbers:interval}{intervalle}).
\end{solution}

\begin{solution}{exo:g12:exp:9}
\emph{1.} Directement de la première inégalité de l'\cref{exo:g12:exp:5}~:
$u_n = \eu^{\,n\ln(1+1/n)} \leq \eu^1 = \eu$.

\emph{2.} Écrire
\[
\ln u_n = n \ln\!\left(1 + \frac1n\right)
= \frac{\ln\!\left(1 + \frac1n\right) - \ln 1}{\frac1n},
\]
un taux d'accroissement de $\ln$ au point $1$ avec incrément
$h = \frac1n \to 0$. Comme $\ln$ est \omterm{def:g12:deriv:derivative}{dérivable} en $1$ de \omterm{def:g12:deriv:derivative}{dérivée} $1$,
$\ln u_n \to 1$. Par \omterm{def:g12:limcont:continuity}{continuité} de $\exp$
(\cref{prop:g12:limcont:seqcont}), $u_n = \eu^{\ln u_n} \to \eu^1 = \eu$.
\end{solution}

\begin{solution}{pb:g12:exp:1}
\textbf{1.} $x = \frac{\ln 5}{2}$. \quad $3x - 1 = \eu^2$, donc
$x = \frac{\eu^2 + 1}{3}$. \quad Avec $u = \eu^x$~:
$u^2 - 3u + 2 = (u - 1)(u - 2) = 0$, donc $x = 0$ ou $x = \ln 2$.

\textbf{2.} $\frac{3\ln 2}{\ln 2} = 3$~; \quad
$3 + \frac12 = \frac72$~; \quad $\frac52$.

\textbf{3.} $f'(x) = \eu^x - 1$~: négative avant $0$, positive
après~: \omterm{def:g10:functions:extrema}{minimum} $f(0) = 0$, donc $f \geq 0$ partout, c'est-à-dire
$\eu^x \geq 1 + x$. En substituant $x = \ln(1 + u)$~:
$1 + u \geq 1 + \ln(1 + u)$, c'est-à-dire $\ln(1 + u) \leq u$.

\textbf{4.} $x^2\eu^{-x} = \frac{x^2}{\eu^x} \to 0$ et
$\frac{\ln x}{x} \to 0$~: les exponentielles écrasent les
puissances, les puissances écrasent les \omterm{def:g12:exp:ln}{logarithmes}. Et
$\frac{\eu^x - 1}{x} = \frac{\eu^x - \eu^0}{x - 0} \to \exp'(0) =
1$.

\textbf{5.} $f'(x) = \ln x + 1$~: nulle en $\eu^{-1}$, négative
avant, positive après~: \omterm{def:g10:functions:extrema}{minimum}
$f\!\left(\frac1\eu\right) = -\frac1\eu$~; et $x \ln x \to 0$ en
$0^+$~: la courbe part de l'origine, plonge jusqu'à $-\frac1\eu$,
puis s'élève.

\textbf{6.} $n = \frac{\ln 2}{\ln 1.03} \approx 23.4$ ans.

\textbf{7.} Valeurs exactes~: $23.4$, $11.9$, $8.0$ ans~; règle de
72~: $24$, $12$, $8$. Comme $\ln(1+r) \approx r$, on a
$\frac{\ln 2}{\ln(1+r)} \approx \frac{0.693}{r}$, c'est-à-dire
$\frac{69.3}{\text{taux en }\%}$~; les banquiers arrondissent à
$72$ parce qu'il se divise magnifiquement par $2, 3, 4, 6, 8, 9,
12$ --- le calcul mental vaut mieux qu'une décimale de précision.

\textbf{8.} $2^{t/20} = \eu^{t \ln 2 / 20}$~:
$\lambda = \frac{\ln 2}{20}$ par minute. Au bout de
$24 \times 60 = 1440$ minutes~: $2^{72} \approx 4.7 \times
10^{21}$ bactéries --- plus que les grains de sable de la Terre, à
partir d'une seule cellule et en un jour. Le modèle n'est honnête
que tant que durent la nourriture et l'espace~: la croissance réelle
s'infléchit en courbe logistique en S (la courbe de
l'épidémiologiste du \cref{pb:g12:deriv:1}).

\textbf{9.} À 23 h (soit $8$ heures plus tard)~:
$100 \times 2^{-8/5} \approx 33$ mg. Pour passer sous $10$ mg~:
$2^{-t/5} < 0.1$ donne $t > 5\,\frac{\ln 10}{\ln 2} \approx
16.6$ h, soit vers 7 h 40 le lendemain matin --- l'expresso de
quinze heures a une longue traîne.

\textbf{10.} Capitalisation annuelle~: $2$. Mensuelle~:
$\left(1 + \frac{1}{12}\right)^{12} \approx 2.613$. Quotidienne~:
environ $2.715$. Le plafond~: $\eu = 2.71828\ldots$
(\cref{exo:g12:exp:9}) --- en capitalisant continûment, la
cupidité \omterm{def:g12:seq:limit}{converge}.

\textbf{11.} Si les cas croissent exponentiellement,
$c(t) = c_0 \eu^{\lambda t}$, alors
$\ln c(t) = \ln c_0 + \lambda t$~: une droite de \omterm{def:g10:reffunc:affine}{coefficient
directeur} $\lambda$ --- le taux de croissance. Une portion
rectiligne sur le \omterm{def:g10:functions:graph}{graphique} logarithmique \emph{est} la phase
exponentielle, et sa pente est le tempo de l'épidémie.

\textbf{12.} $120 - 60 = 60$ dB signifie
$10\log_{10}(I_2/I_1) = 60$~: le concert est $10^6$ fois --- un
million de fois --- plus intense que la conversation.

\textbf{13.} Amplitude~: $10^{7-5} = 100$. Énergie~:
$100^{3/2} = 1000$~: deux points de magnitude cachent trois ordres
de grandeur en énergie.

\textbf{14.} $10^{6.5 - 2} = 10^{4.5} \approx 32\,000$~: le jus de
citron est trente mille fois plus acide que le lait --- le pH
comprime les abîmes de la chimie sur une échelle de poche.

\textbf{15.} $\log_2\!\frac{4186}{27.5} =
\frac{\ln(4186/27.5)}{\ln 2} \approx 7.25$~: un piano couvre un peu
plus de sept octaves. Nos sens réagissent aux \emph{rapports} de
stimulus --- doubler l'intensité se ressent comme un pas, quel que
soit le niveau de départ --- si bien que la perception est bâtie sur
une échelle logarithmique, comme les unités que nous avons inventées
pour elle.

\textbf{16.} Placer la réglette de $a$ bout à bout avec la position
de $b$ additionne les longueurs $\ln a + \ln b$, et la graduation
située à cette longueur totale se lit
$\eu^{\ln a + \ln b} = ab$~: la règle à calcul calcule des produits
en additionnant des \omterm{def:g12:exp:ln}{logarithmes} --- $\ln(ab) = \ln a + \ln b$
(\cref{prop:g12:exp:lnalg}) gravé dans le buis.

\textbf{17.} Le \omterm{def:g10:functions:extrema}{minimum} de $\frac{\eu^x}{x}$ sur
$\intoo{0}{+\infty}$ vaut $\eu$, atteint en $x = 1$~: aucune
solution pour $k < \eu$, exactement une pour $k = \eu$, deux pour
$k > \eu$. Vérification de la \omterm{def:g12:deriv:tangent}{tangente}~: en $a = 1$, la \omterm{def:g12:deriv:tangent}{tangente} à
$\eu^x$ est $y = \eu + \eu(x - 1) = \eu x$~: elle passe par
l'origine --- la droite seuil, qui effleure la courbe en
$(1, \eu)$.

\textbf{18.} $n \ln\!\left(1 - \frac1n\right) =
\frac{\ln(1 - 1/n)}{-1/n} \times (-1) \to -1$, donc
$\left(1 - \frac1n\right)^n \to \eu^{-1} \approx 0.368$~: le
$0.999^{1000} \approx 0.368$ du \cref{pb:g11:binom:1}, enfin
expliqué --- la constante du presque-raté de la loterie est
$\frac1\eu$.

\textbf{19.} Les deux sont la forme limite de «~beaucoup
d'\omterm{def:g11:prob:model}{événements} indépendants, chacun individuellement improbable~»~: la
\omterm{def:g10:proba:distribution}{probabilité} qu'\emph{aucune} des $n$ chances de taille $\frac1n$ ne
se réalise tend vers $\eu^{-1}$, et la règle du recrutement règle sa
fenêtre de rejet de sorte que le succès se concentre sur cette même
constante. $\frac1\eu \approx 0.368$.

\textbf{20.} Le plafond de la capitalisation~; l'unique base dont la
\omterm{def:g12:deriv:tangent}{tangente} en $0$ a un \omterm{def:g10:reffunc:affine}{coefficient directeur} égal à $1$ (raison pour
laquelle l'analyse l'adore)~; la croissance qui distance toutes les
puissances~; la constante où se rangent les presque-ratés et l'arrêt
optimal~; et son inverse, le \omterm{def:g12:exp:ln}{logarithme} --- la multiplication
devenue addition, les plus folles amplitudes du monde repliées sur
une règle graduée.
\end{solution}
